\section{引言}
\label{sec:introduction}

对部分子分布（PDFs）的精确理解~\cite{Forte:2010dt,Forte:2013wc,Ball:2015oha}在希格斯玻色子的发现中发挥了重要作用，并将是 LHC 上寻找新物理的关键要素~\cite{Rojo:2015acz}。近年来，一批新一代 PDF 集合~\cite{Ball:2014uwa,Dulat:2015mca,Harland-Lang:2014zoa,Alekhin:2017kpj,Abramowicz:2015mha,Jimenez-Delgado:2014twa,Accardi:2016qay}已被发展出来供 LHC Run II 使用。其中若干集合被用于构造 PDF4LHC15 组合集合，该组合被推荐用于新物理搜寻以及精密观测量 PDF 不确定度的评估~\cite{Butterworth:2015oua}。这些 PDF4LHC15 集合是通过对三个全局拟合~\cite{Ball:2014uwa,Dulat:2015mca,Harland-Lang:2014zoa}进行统计组合得到的：其正当性在于各全局确定之间的符合程度已大为改善，它们之间的差异大体上与统计涨落相一致。

尽管有这些进展，PDF 确定在精度与可靠性方面仍需改进。LHC 上的精密测量，例如通过希格斯耦合测量对新物理的搜寻，最终将要求百分之一量级的系统性 PDF 知识，以充分发挥 LHC 的潜力。NNPDF3.0 PDF 集合~\cite{Ball:2014uwa}是进入 PDF4LHC15 组合的集合之一，其独特之处在于它基于一套经过封闭性检验系统验证的方法学，这类检验保证了从数据提取 PDF 的流程在统计上的自洽性。本文的目标是给出 NNPDF3.0 集合的更新版本 NNPDF3.1，这是迈向百分之一水平不确定度 PDF 的第一步。要达到这一目标需要在两个方向上取得进展，因此更新的动机也是双重的。

一方面，把 PDF 精度降到百分之一水平需要更大、更精确的数据集，以及相应精确的理论预言。自 NNPDF3.0 发布以来，已有大量新的实验测量结果可供使用。来自 Tevatron 的有：基于完整 Run II 数据集、以电子和 $\mu$ 子末态测量的 $W$ 玻色子不对称度最终结果~\cite{Abazov:2013rja,D0:2014kma}。在 LHC 方面，ATLAS、CMS 与 LHCb 实验发布了对单举喷注产生、规范玻色子产生和顶夸克产生的多种多样测量。最后，HERA 的 DIS 结构函数合并终版测量也已可用~\cite{Abramowicz:2015mha}。与实验进展并行的是，近期完成了数量可观的对 PDF 有直接敏感性的强子对撞过程高精度 QCD 计算，使得这些过程可用于 NNLO 水平的 PDF 确定。其中包括顶夸克对产生的微分分布~\cite{Czakon:2015owf,Czakon:2016dgf}、$Z$ 与 $W$ 玻色子的横动量~\cite{Boughezal:2015dva,Ridder:2015dxa}以及单举喷注产生~\cite{Currie:2013dwa,Currie:2016bfm}，对它们都有 ATLAS 与 CMS 的精密数据集可用。

所有这些新数据集与计算都已被纳入 NNPDF3.1。新数据的纳入带来了新的挑战。鉴于其中一些测量所基于的数据集非常庞大，非关联实验不确定度往往达到千分之一量级。要实现良好拟合就需要对关联系统误差和理论预言的数值精度都有空前的控制能力。

另一方面，当不确定度达到百分之一量级时，此前未纳入 PDF 确定的理论不确定度相关的精度问题就变得重要起来。虽然在 PDF 确定中全面纳入理论不确定度尚需进一步研究，我们近期已经论证：一个重要的理论偏差来源来自“粲完全由胶子和轻夸克微扰生成”这一传统假设。一种允许在 FONLL 匹配的一般质量变味数方案中纳入参数化重夸克 PDF 的方法学已被发展出来~\cite{Ball:2015tna,Ball:2015dpa}，并在一次 NNPDF PDF 确定中得到实现~\cite{Ball:2016neh}。研究发现，当粲 PDF 与其他 PDF 一起由数据参数化并确定时，与粲质量取值相关的不确定度大部分转化为标准的 PDF 不确定度，而与“粲 PDF 纯粹是微扰的”这一假设相关的偏差则被消除~\cite{Ball:2016neh}。因此在 NNPDF3.1 中粲被作为独立 PDF 参数化，与轻夸克和胶子同等对待。我们将证明这在不增加不确定度的情况下改善了拟合质量，并稳定了 PDF 对粲质量的依赖，对轻夸克 PDF 而言这种依赖几乎被完全消除。

NNPDF3.1 PDF 集合以 LO、NLO 和 NNLO 精度发布。首次地，全部 NLO 与 NNLO PDF 同时以 Hessian 集合和蒙特卡罗副本两种形式提供，并利用了近期发展的构造最优 Hessian 表示的强大方法~\cite{Carrazza:2015aoa}。此外，同样首次地，默认 PDF 集合以压缩蒙特卡罗集合的形式提供~\cite{Carrazza:2015hva}。因此尽管只呈现为 100 个蒙特卡罗副本的集合，它们却展现出大得多的集合的许多统计性质，在不损失信息的前提下减少了观测量计算时间。借助 {\tt SM-PDF} 工具~\cite{Carrazza:2016htc}还可进一步提高计算效率：该工具可为特定过程或过程类的不确定度计算选取最优的 Hessian 本征矢量子集，并已作为网页界面~\cite{Carrazza:2016wte}提供，现已包含 NNPDF3.1 集合。我们还提供了多种基于数据子集的 PDF 集合（标准形式为 100 副本蒙特卡罗集合），它们对诸如新物理搜寻或标准模型参数测量等特定应用可能有用。

本文结构如下。
%
首先，在第~\ref{sec:expdata} 节中我们讨论新数据集的实验方面及相关理论问题。
随后在第~\ref{sec:results} 节中详细描述基线 NNPDF3.1 PDF 集合，并专门讨论方法学改进的影响，特别是粲 PDF 如今像其他 PDF 一样被独立参数化和确定。在第~\ref{sec:impactnewdata} 节中，我们通过比较基于不同数据子集的 PDF 集合讨论新数据的影响，并讨论基于更保守数据子集的 PDF 集合。在第~\ref{sec:pheno} 节中我们总结 PDF 与亮度的不确定度状况，并结合我们的结果专门讨论质子的奇异味与粲含量，同时给出 LHC 的初步唯象学研究。最后，第~\ref{sec:conclusion} 节总结了以各种格式发布的 PDF 集合，并给出可下载更详细图表的存储库链接。

%%%%%%%%%%%%%%%%%%%%%%%%%%%%%%%%%%%%%%%%%%%%%%%%%%%%%%%%%%%%%%%%%%%%%%%%%%%
%%%%%%%%%%%%%%%%%%%%%%%%%%%%%%%%%%%%%%%%%%%%%%%%%%%%%%%%%%%%%%%%%%%%%%%%%%%
